% ============================================================
% PACKAGES
% ============================================================

\usepackage[T1]{fontenc}
\usepackage[utf8]{inputenc}
\usepackage[ngerman]{babel}
\usepackage[sfdefault,lf]{carlito} % Carlito ist in LaTeX das Äquivalent zu Calibri (Open-Source von Google), 'lf' erzwingt Lining Figures (normale Zahlenhöhe) statt Old-Style (Zahlen wie ein kleines o)
\usepackage{sansmath}          % Mathe-Modus auch in Sans-Serif
\sansmath
\linespread{0.9} % Etwas engere Zeilen für mehr Kompaktheit bei 10pt

\usepackage[landscape, a4paper, margin=1.5mm]{geometry}
\usepackage{microtype}
\usepackage{multicol}

\usepackage{mathtools}          % Superset von amsmath: DeclarePairedDelimiter, coloneqq, etc.
\let\Bbbk\relax % Verhindert Konflikt zwischen newpxmath und amssymb
\usepackage{amssymb}
\usepackage{siunitx}
\sisetup{
  mode=match,
  propagate-math-font=true,
  reset-text-family=false,
  exponent-product = \cdot,
  output-decimal-marker = {,},
  per-mode = symbol
}

\usepackage[table]{xcolor}
\usepackage{array}
\usepackage{tabularx}
\usepackage{tabularray}
\usepackage[column=O]{cellspace}
\usepackage{ragged2e} % Erlaubt Silbentrennung bei linksbündigem/zentriertem Flattersatz in schmalen Spalten
\usepackage{xstring}
\usepackage{varwidth}

\usepackage{graphicx}
\usepackage{tikz}
\usepackage{pgfplots}
\pgfplotsset{compat=1.18}
\usepackage{tcolorbox}
\tcbuselibrary{skins, breakable, raster}

\usepackage{hyperref}
\usepackage{enumitem}
\usepackage{makeidx}               % Stichwortverzeichnis (Setup in styles/65_index)
